\thispagestyle{empty}
\begin{center}
\vspace*{1.2cm}
{\sffamily\bfseries\Huge Cálculo II}\\[6pt]
{\sffamily\bfseries\LARGE Resumo para a Prova 1}\\[18pt]
{\sffamily\large Somas de Riemann \;\textperiodcentered\; Integrais definidas e indefinidas \;\textperiodcentered\; Como \emph{ler} uma derivada}\\[8pt]
{\small Baseado nas questões recomendadas pela professora --- Stewart, \emph{Cálculo}, vol.~1, 9\textsuperscript{a} ed., seções 4.9 a 6.3}\\[4pt]
{\small Ciência da Computação --- 2\textsuperscript{o} período --- PUC Minas --- \today}
\end{center}

\vspace{0.6cm}

\begin{tcolorbox}[base, colback=azulclaro, colframe=azul, title={Como usar este resumo}]
\begin{itemize}[leftmargin=1.4em, itemsep=2pt]
\item \textbf{Parte 1} é a base: se você lê corretamente o que uma derivada \emph{significa}, toda integral vira ``desfazer uma derivada'' e todo teorema do capítulo 5 vira óbvio.
\item \textbf{Parte 3} (somas de Riemann) é o coração da prova: cada símbolo da fórmula é explicado, com um exemplo completo calculado de três formas diferentes (soma com $n=6$, limite exato, TFC) --- e as três batem.
\item Cada exemplo resolvido segue o padrão: \emph{classificar} $\to$ \emph{resolver sem pular passos} $\to$ \emph{interpretar geometricamente} $\to$ \emph{verificar}. Depois de cada seção há exercícios similares; o gabarito está no fim.
\item As caixas \textcolor{vermelho}{\textbf{vermelhas}} são os erros que custam ponto. Leia-as duas vezes.
\end{itemize}
\end{tcolorbox}

\tableofcontents
\newpage
